%---------------------------------------------------------------------
\chapter{Observability and Evaluation}
%---------------------------------------------------------------------

This chapter describes what is recorded for every run and reports how well the agent performs. All
numbers come from runs of the scripts \texttt{loomtrace/experiments.py} and \texttt{run\_tests.py} on
11 November 2026 on the hardware described in Section~3.5; the raw output is stored in
\texttt{results/experiments.json} and \texttt{TEST\_RESULTS.md}. All runs used the deterministic rule
mode, without a language model.

\section{Tracing and Monitoring}

Three levels of information are captured for every case.

\begin{enumerate}[itemsep=2pt]
\item \textbf{Node trace.} Each node appends an entry with its name, a summary message and a
timestamp; the planner's entry includes the plan and query, the verifier's includes every check with
its status, reason and source, and the grader's includes completeness, consistency, citation fraction
and its decision. Listings~\ref{lst:replan} and~\ref{lst:hitl} show two such traces.
\item \textbf{Tool log.} Each tool call is recorded with its iteration, input-derived source
identifier, success flag, note and elapsed time in milliseconds (Listing~\ref{lst:tools}).
\item \textbf{Checkpoints.} After every node, the complete state is serialised to JSON and stored in the
\texttt{checkpoints} table with the case identifier, step number and node name
(Listing~\ref{lst:ckpt}). The \texttt{cases}, \texttt{review\_queue} and \texttt{decisions} tables
record the outcome of each case and of each human review.
\end{enumerate}

Because all model calls were disabled, token usage and monetary cost were zero for every run and are
not reported further. Latency was measured with a high-resolution timer around each complete graph
invocation, including the SQLite writes. Over three repetitions of the 45-case set, the mean latency
of the full agent was 16.7\,ms per case (standard deviation of the three repetition means 2.9\,ms).
The slowest single case took 391.5\,ms; it was the first case of the first repetition and includes
one-time start-up work. In the second and third repetitions, no case took longer than 25\,ms. These values are four
orders of magnitude below the 15-second requirement (NFR1), as expected for a configuration without
model calls.

\section{Evaluation Dataset}

The golden set contains 45 cases: 3 products (Kancheepuram Silk, Pochampalli Ikat, and Banaras
Brocades and Sarees) $\times$ 3 verdict classes $\times$ 5 cases, the structure specified in the design
document. Each case is a label text with product, region and state, and, depending on the case,
material, weaver identifier and batch identifier. The cases are generated by the script
\texttt{evaluate.py} from the registry and the simulated producer records according to fixed rules:

\begin{itemize}[itemsep=2pt]
\item \textbf{Supported (5 per product).} A real weaver of the product and one of that weaver's
batches, the weaver's district as region, and the batch's yarn as material. Two of the five use a
market alias instead of the registered name (``Kanjivaram Silk Saree'', ``Pochampally Ikat Saree'',
``Banarasi Silk Saree'').
\item \textbf{Conflicting evidence (5 per product).} One case for each of five contradictions: a region
outside the GI area (Surat, Bhagalpur or Murshidabad); a weaver and batch from a different product's
cluster; a batch belonging to a different weaver of the same product; a product name that is not a
covered GI (``Mysore Silk Saree''); and a material inconsistent with the registration (polyester).
\item \textbf{Insufficient evidence (5 per product).} Weaver missing; batch missing; both missing; an
unknown weaver identifier (\texttt{WV-XXX-999}); and an unknown batch identifier (\texttt{B-XXX-0000}).
\end{itemize}

The label of each case is assigned by construction from these definitions, which were written by the
team before evaluation; the labels were not assigned independently by two annotators. The class
distribution is balanced, 15 cases per class. The test set was not fully separate from development:
as described in Section~7.6, one verification rule was changed after the first evaluation run revealed
three errors on this set. In addition, a separate set of 15 hand-written demonstration cases
(\texttt{run\_tests.py}) covers nine of the fifteen products, including four products that are not in
the golden set (Sambalpuri Bandha, Muga Silk, Kota Doria and Chanderi Sarees) and Kullu Shawls.

\section{Metrics and Evaluation Rubric}

Let $N = 45$ and let $y_k$ and $\hat{y}_k$ be the labelled and predicted verdicts of case $k$. The
following metrics are used.

\begin{itemize}[itemsep=2pt]
\item \textbf{Verdict accuracy} $= \frac{1}{N}\sum_k \mathbb{1}[\hat{y}_k = y_k]$.
\item \textbf{Per-class precision, recall and F1}, with $F_1 = \frac{2PR}{P+R}$, and \textbf{macro-F1}, the
unweighted mean of the three class F1-scores.
\item \textbf{Correct abstention}: the fraction of cases labelled \emph{Insufficient evidence} that are
predicted as \emph{Insufficient evidence} and therefore sent to human review.
\item \textbf{False Supported}: the number of cases predicted \emph{Supported} whose label is not
\emph{Supported}. This is the most harmful error, because it endorses an unsupported claim.
\item \textbf{Retrieval Recall@5}: the fraction of cases in which the final document search returns at
least one passage of the correct product among its top five results.
\item \textbf{Citation rate}: the fraction of decided verdicts (\emph{Supported} or \emph{Conflicting})
that list at least one source identifier.
\item \textbf{Latency and tool calls}: mean wall-clock time and mean number of tool calls per case.
\end{itemize}

The design document also proposed \emph{citation faithfulness}, the share of verdict statements backed
by the cited text, measured manually. This was not measured: the citation rate above records only that
citations are present, not that each cited passage supports the statement. No model was used as a
judge, because every output is a discrete label or check status with an exact expected value, for
which exact-match metrics are more reliable than a model judge~\cite{zheng2023judge}.

\section{Baselines}

Two comparison systems were run on the same 45 cases and the same data.

\begin{enumerate}[itemsep=2pt]
\item \textbf{Single-lookup baseline.} This follows the baseline of the design document: one document
search using the whole label as the query, followed by one decision. Because no language model was
used, the decision step is the rule ``\emph{Supported} if the top passage scores above 2, otherwise
\emph{Insufficient evidence}''. The baseline has no planning, no producer checks, no loop and no
guardrail. It is a weak baseline, and it is reported to show what retrieval alone achieves.
\item \textbf{LoomTrace without re-planning.} The full agent with the iteration limit set to one, so
that the grader can never loop back. This is the variant~A of the experiment in Chapter~8 and isolates
the contribution of the agentic loop.
\end{enumerate}

\section{Results}

Table~\ref{tab:results} summarises the results. Because the system and the baseline are deterministic,
the three repetitions produced identical verdicts, and the standard deviation of every quality metric
is zero; only latency varied between repetitions.

\begin{table}[H]
\centering
\caption{Evaluation results on the 45-case golden set (mean of 3 repetitions)}
\label{tab:results}
\small
\begin{tabularx}{\linewidth}{@{}>{\raggedright\arraybackslash}X c c c@{}}
\toprule
\textbf{Metric} & \textbf{Single lookup} & \textbf{LoomTrace, no re-plan} & \textbf{LoomTrace (full)} \\ \midrule
Verdict accuracy & 33.3\% & 80.0\% & \textbf{100.0\%} \\
Macro-F1 & 0.17 & 0.80 & \textbf{1.00} \\
F1 Supported / Conflicting / Insufficient & 0.50 / 0.00 / 0.00 & 0.75 / 0.89 / 0.77 & 1.00 / 1.00 / 1.00 \\
Correct abstention (15 cases) & 0\% & 100\% & 100\% \\
False Supported & 30 & 0 & 0 \\
Retrieval Recall@5 & -- & 1.00 & 1.00 \\
Citation rate (decided verdicts) & -- & 100\% & 100\% \\
Mean citations per case & -- & 6.40 & 6.53 \\
Mean tool calls per case (max) & 1 (1) & 6.27 (7) & 7.07 (11) \\
Cases that re-planned & -- & 0 & 9 \\
Cases escalated to review & -- & 24 & 15 \\
Mean latency per case & -- & 13.7\,ms & 16.7\,ms \\
\bottomrule
\end{tabularx}
\end{table}

As Table~\ref{tab:results} shows, the single-lookup baseline labels every case \emph{Supported}: the
label always contains the product name, so a passage about the product is always retrieved, and the
baseline has no means of detecting a contradiction or a missing identifier. It therefore endorses all
30 conflicting and insufficient claims. The full agent classifies all 45 cases correctly, sends exactly
the 15 insufficient cases to human review, and never issues a false \emph{Supported} verdict. Nine cases
used the re-planning loop: the six \emph{Supported} cases with alias names and the three
\emph{Conflicting} cases with an unregistered product name. Table~\ref{tab:confusion} gives the
confusion matrices.

\begin{table}[H]
\centering
\caption{Confusion matrices (rows: label; columns: prediction; S = Supported, C = Conflicting, I = Insufficient)}
\label{tab:confusion}
\small
\begin{tabular}{@{}l ccc c ccc c ccc@{}}
\toprule
& \multicolumn{3}{c}{\textbf{Single lookup}} && \multicolumn{3}{c}{\textbf{No re-plan}} && \multicolumn{3}{c}{\textbf{Full}} \\
\cmidrule{2-4}\cmidrule{6-8}\cmidrule{10-12}
\textbf{Label} & S & C & I && S & C & I && S & C & I \\ \midrule
Supported & 15 & 0 & 0 && 9 & 0 & 6 && 15 & 0 & 0 \\
Conflicting & 15 & 0 & 0 && 0 & 12 & 3 && 0 & 15 & 0 \\
Insufficient & 15 & 0 & 0 && 0 & 0 & 15 && 0 & 0 & 15 \\
\bottomrule
\end{tabular}
\end{table}

Figure~\ref{fig:f1} compares the per-class F1-scores of the three systems.

\begin{figure}[H]
\centering
\begin{tikzpicture}
\begin{axis}[ybar,bar width=11pt,width=0.9\linewidth,height=6cm,
  ymin=0,ymax=1.15,ylabel={F1-score},
  symbolic x coords={Supported,Conflicting,Insufficient},xtick=data,
  nodes near coords,nodes near coords align={vertical},
  every node near coord/.append style={font=\scriptsize,/pgf/number format/fixed,/pgf/number format/precision=2},
  enlarge x limits=0.25,axis lines*=left,ymajorgrids,grid style={gray!25},
  legend style={at={(0.5,-0.18)},anchor=north,legend columns=3,font=\small,draw=none}]
\addplot[fill=gray!40,draw=gray] coordinates {(Supported,0.50) (Conflicting,0.00) (Insufficient,0.00)};
\addplot[fill=bmsblue!35,draw=bmsblue] coordinates {(Supported,0.75) (Conflicting,0.889) (Insufficient,0.769)};
\addplot[fill=bmsblue!85,draw=bmsblue] coordinates {(Supported,1.00) (Conflicting,1.00) (Insufficient,1.00)};
\legend{Single lookup, No re-plan, Full LoomTrace}
\end{axis}
\end{tikzpicture}
\caption{Per-class F1-score on the 45-case golden set}
\label{fig:f1}
\end{figure}

Table~\ref{tab:targets} compares the full agent with the targets set in the design document.

\begin{table}[H]
\centering
\caption{Results against the targets of the design document}
\label{tab:targets}
\small
\begin{tabularx}{\linewidth}{@{}>{\raggedright\arraybackslash}X c c >{\raggedright\arraybackslash}p{3.6cm}@{}}
\toprule
\textbf{Metric} & \textbf{Target} & \textbf{Result} & \textbf{Status} \\ \midrule
Verdict accuracy & $\geq 80\%$ & 100\% & Met \\
Retrieval Recall@5 & $\geq 0.80$ & 1.00 & Met \\
Citation faithfulness & $\geq 90\%$ & not measured & Not assessed; citation presence 100\% \\
Correct abstention & $\geq 90\%$ & 100\% & Met \\
Latency & $< 15$\,s & 16.7\,ms & Met (rule mode only) \\
\bottomrule
\end{tabularx}
\end{table}

On the 15 demonstration cases, all 15 verdicts matched the expected verdict: 7 \emph{Supported} (including
Sambalpuri Bandha, Muga Silk and Kota Doria), 5 \emph{Conflicting} (wrong region, foreign weaver, batch
of another weaver, polyester material, and a Chanderi saree claimed from Varanasi) and 3
\emph{Insufficient} (missing identifiers, unknown weaver, unknown batch). These cases are also
team-written and therefore provide functional coverage of the additional products rather than an
independent measure of accuracy.

\section{Error Analysis}

The full agent made no errors on the final golden set. This result must be read with the construction of
the test set in mind, and the analysis below therefore covers four sources of errors: errors observed
during development, errors of the ablated agent and the baseline, and errors on realistic inputs that
the golden set does not contain.

\paragraph{Errors found and corrected during development.} The first evaluation run of the full agent
reached 93.3\% accuracy (42 of 45) and 80\% correct abstention. All three errors were cases with an
unknown weaver identifier and a valid batch, labelled \emph{Insufficient} but predicted
\emph{Conflicting}, because the batch's weaver did not match the unknown identifier. The rule was
changed so that a batch--weaver mismatch is reported as \emph{unknown} when the claimed weaver itself
cannot be verified; after this change all 45 cases were correct. Because the change was made after
observing the test set, the final 100\% is partly a result of fitting to it. Separately, the
demonstration tests revealed that a claimed material of ``silk'' was rejected for Kancheepuram Silk,
whose registered material is ``mulberry silk''; the material comparison was changed to accept a match in
either direction.

\paragraph{Errors without re-planning.} All nine errors of the ablated agent fall into two categories:
six alias names that were not resolved by exact matching (\emph{Supported} predicted as
\emph{Insufficient}) and three unregistered product names that could not be classified as
\emph{Conflicting} on the first iteration. These are exactly the nine cases that the full agent resolves
by re-planning, which is analysed in Chapter~8.

\paragraph{Errors of the baseline.} All 30 errors are false \emph{Supported} verdicts, since the baseline
never checks region, state, material or producer records.

\paragraph{Errors on realistic inputs outside the test set.} To estimate behaviour beyond the constructed
cases, six further inputs that vary the spelling and form of the label were run through the full agent.
Table~\ref{tab:probes} shows the results. Four of the six produced a wrong \emph{Conflicting} verdict.

\begin{table}[H]
\centering
\caption{Behaviour on six realistic inputs not in the golden set}
\label{tab:probes}
\small
\begin{tabularx}{\linewidth}{@{}>{\raggedright\arraybackslash}p{4.4cm} >{\raggedright\arraybackslash}p{2.3cm} >{\raggedright\arraybackslash}X@{}}
\toprule
\textbf{Input variation} & \textbf{Verdict} & \textbf{Cause} \\ \midrule
Alias not in list (``Conjeevaram Silk Saree'') & Conflicting (wrong) & Product unresolved after re-plan, so the product check fails \\
District spelling (``Kanchipuram'') & Conflicting (wrong) & Region compared by substring; spelling variant not matched \\
State abbreviation (``TS'') & Conflicting (wrong) & State compared by exact name; census lookup also fails \\
Former state name (``Andhra Pradesh'' for Pochampalli) & Conflicting (wrong) & Registry stores current state (Telangana) only \\
Lower-case identifiers & Supported (correct) & Identifiers are upper-cased before lookup \\
Free-text sentence with alias and identifiers & Supported (correct) & Parser finds product alias and identifier patterns \\
\bottomrule
\end{tabularx}
\end{table}

These four failures share one cause: the region, state and product comparisons are exact or substring
matches against a single canonical spelling. They would be fixed by (i) fuzzy matching of district and
product names against the registry, with a similarity threshold below which the check is
\emph{unknown} rather than \emph{fail}; (ii) a table of state abbreviations and former state names; and
(iii) treating a product that cannot be resolved as ``not covered'' inside the agent, as the interface
already does before the agent runs. The more general lesson is that a contradiction should only be
reported when the mismatch is certain, and that uncertain mismatches should be routed to human review.
